% ECONOMICS_PROBLEM: {"id": "C72-153", "title": "Exact Nash complexity in concisely represented potential games", "jel_code": "C72", "source_id": "OP-0139"}
% !TeX program = xelatex
\documentclass[11pt,a4paper]{article}
\usepackage[margin=23mm]{geometry}
\usepackage{fontspec}
\setmainfont{Latin Modern Roman}
\usepackage{amsmath,amssymb,mathtools}
\usepackage{xcolor,enumitem,booktabs,longtable,array,xurl,hyperref}
\definecolor{muted}{HTML}{53636C}
\hypersetup{colorlinks=true,urlcolor=blue,linkcolor=blue}
\setlength{\parindent}{0pt}
\setlength{\parskip}{5pt}
\newcommand{\R}{\mathbb R}
\newcommand{\E}{\mathbb E}
\newcommand{\N}{\mathbb N}
\newcommand{\ind}{\mathbf 1}
\newcommand{\Var}{\operatorname{Var}}
\newcommand{\Cov}{\operatorname{Cov}}
\newcommand{\supp}{\operatorname{supp}}
\DeclareMathOperator*{\argmin}{arg\,min}
\DeclareMathOperator*{\argmax}{arg\,max}
\newcommand{\entrymeta}[3]{{\sffamily\small\color{muted}\textbf{\texttt{#1}}\quad|\quad #2\par #3\par}}
\newcommand{\Problemref}[1]{\texttt{#1}}
\newcommand{\JELref}[2]{\texttt{#2} (JEL \texttt{#1})}
\begin{document}
\section*{Exact Nash complexity in concisely represented potential games}
\textbf{Problem ID:} \texttt{C72-153}\quad \textbf{JEL:} \texttt{C72}\par
% BEGIN ECONOMICS STATEMENT
\label{problem:OP-0139}
% GAME_THEORY_PROBLEM: {"local_id": "GT-009", "original_number": 9, "kind": "P", "entry_kind": "statement_only", "published": false, "review_required": true, "category": "game_theory"}
\label{game:GT-009}
{\sffamily\small\color{muted}
\textbf{GT-009} | Potential games and search complexity\\
\textbf{P} | Explicit source question; succinct representation is essential.
\par}
\subsection*{Definitions and mathematical statement}
Let $S_i$ be explicitly listed finite action sets and $X=\prod_{i=1}^k\Delta(S_i)$; $k$ is part of the input. An exact potential game admits $\Phi:\prod_iS_i\to\mathbb Q$ such that
\[
 u_i(a_i,a_{-i})-u_i(b_i,a_{-i})
 =\Phi(a_i,a_{-i})-\Phi(b_i,a_{-i})
 \quad\text{for all }i,a_i,b_i,a_{-i}.
\]
The input provides the multilinear extension $\Phi:X\to\mathbb R$ by a succinct rational arithmetic circuit. The source's circuit convention permits addition and multiplication; inputs to a multiplication gate involve disjoint player blocks and disjoint subcircuits, so exact rational evaluation has polynomial bit cost. Payoff differences, and hence Nash equilibria, are determined by $\Phi$.
The exact search target is any real $x\in X$ satisfying
\[
 \Phi(x)\geq\Phi(p_i,x_{-i})\quad
 (i\in\{1,\ldots,k\},\ p_i\in\Delta(S_i)).
\]
Under the real-search reductions and PLS convention of [S1], is this search problem complete for $\mathsf{FIXP}\cap\mathsf{PLS}$?
\subsection*{Short English statement}
What is the precise complexity of computing an exact equilibrium when a potential is available through a compact circuit?
\subsection*{Variants and scope boundaries}
``Mixed strategy'' allows pure profiles. With a fully expanded payoff table, scanning pure profiles for a potential maximum is polynomial in that table's size, so the succinct encoding cannot be omitted. $\mathsf{FIXP}$ uses exact fixed points of rational algebraic circuits; it does not demand a finite decimal representation of an arbitrary irrational point. $\mathsf{PLS}$ describes reduction to finite local search with polynomial-time neighborhood and value computations.
\subsection*{Source relationship and status}
Question 6.1 states the completeness question. Its footnote explicitly specifies a PLS convention that does not require efficient verification of every purported target solution. A different convention must not silently be substituted.
\subsection*{References}
\begingroup\footnotesize\raggedright
{\footnotesize\raggedright
\textbf{[S1]} Ioannis Anagnostides, Maria-Florina Balcan, Kiriaki Fragkia, Tuomas Sandholm, Emanuel Tewolde, and Brian Hu Zhang (2026 version). \emph{The Complexity of Equilibrium Refinements in Potential Games}. arXiv:2511.03968, PDF dated 11 February 2026. Locator: Definitions 3.3--3.4; Section 3.1; Question 6.1 and footnote 5, printed p.\ 20. \url{https://arxiv.org/pdf/2511.03968}\par}
\par\endgroup

\subsection*{Common conventions: Source mathematical conventions}

\textbf{Finite games.} For a finite set $A$, $\Delta(A)$ is its probability
simplex. Players choose independent mixed strategies in Nash equilibrium;
correlation is allowed only when explicitly part of the solution concept.
In a game with payoff functions $u_i$ and strategy profile $x$, an additive
$\varepsilon$ Nash equilibrium satisfies
\[
 u_i(x)\geq u_i(x_i',x_{-i})-\varepsilon
 \quad\text{for every player $i$ and every admissible deviation $x_i'$}.
\]
For numerical approximation, payoffs are normalized as locally specified.
An $\varepsilon$ payoff residual is distinct from distance at most
$\varepsilon$ to the set of exact equilibria. Point convergence, convergence
of empirical play, and convergence in distance to an equilibrium set are
also distinct.

\textbf{Computational representations.} Unless stated otherwise, a complexity
question takes rational data with binary encoding; polynomial time is measured
in total bit length. A fixed-accuracy polynomial algorithm is not necessarily
polynomial in $1/\varepsilon$, and neither is necessarily polynomial in
$\log(1/\varepsilon)$. Succinct games and value, demand, or sampling oracles
must declare their interfaces. Exact real solutions require an explicit
representation; an unspecified list of arbitrary real numbers is not a
finite computational output. A claimed algorithmic barrier is meaningful
only for its specified model and reductions.

\textbf{Dynamic games.} Histories, information, observation, and allowable
randomization are part of the model. A stationary strategy depends only on
the current state; a positional strategy in a deterministic graph game
chooses one action at each controlled vertex. Nash equilibrium at an initial
state and equilibrium after every history are different requirements.
An expectation of a pathwise $\liminf$ payoff, a $\liminf$ of expected averages,
a limiting expected average, and a uniform finite-horizon equilibrium are
different criteria. None is silently substituted for another.

\textbf{Allocation and mechanisms.} A complete allocation assigns every item
exactly once unless otherwise stated. Zero-valued goods, negative values,
capacity constraints, feasibility, lotteries, and copies of goods affect
fairness definitions and are specified locally. Dominant-strategy incentive
compatibility, Bayesian incentive compatibility, universal truthfulness,
and truthfulness in expectation impose different inequalities. Whenever
an objective ratio has zero denominator, its convention is stated or those
instances are excluded.

\textbf{Infinite-dimensional models.} Expectations, integrals, stochastic
equations, and optimization objectives require their local regularity,
integrability, filtrations, and admissible domains. A backward--forward
mean-field system is a boundary-value problem; it is not an initial-value
semiflow without an additional construction. Long-time, large-population,
and vanishing-noise limits require an order of limits and a convergence
topology. If a minimum need not be attained, the objective is its infimum
or supremum and the approximation requirement is stated separately.
% END ECONOMICS STATEMENT
\end{document}
